\documentclass[12pt]{article}
\usepackage[a4paper, margin=1in]{geometry}
\usepackage{amsmath, amssymb}
\usepackage{booktabs}

\title{\textbf{LECTURE 23 SCRIBE}\\
Predicting the Impossible: Tail Bounds and System Reliability}
\author{}
\date{}

\begin{document}

\maketitle

\section*{Case Study: Amazon Black Friday Sale 2026}

\section*{Slide 1: Lecture Title}
\textbf{Instructor:}
\begin{itemize}
    \item Dhaval Patel, PhD
    \item Associate Professor, CSE
    \item MICxN Research Lab
    \item SEAS, Ahmedabad University
\end{itemize}

\section*{Slide 2: Amazon Black Friday Sale 2026}
\textbf{Details:}
\begin{itemize}
    \item Expected Start Date: 27 November 2026 (Friday)
    \item Based on historical patterns and forecasts
\end{itemize}

\textbf{Observation:}
\begin{itemize}
    \item Large-scale uncertainty
    \item Heavy traffic and unpredictable user behavior
\end{itemize}

\section*{Slide 3: Uncertainty in Computing Systems}
\begin{itemize}
    \item Real-world systems operate under uncertainty
    \item User requests and traffic are random
\end{itemize}

\textbf{Challenge:} Exact prediction is not possible.

\section*{Slide 4: Motivation}
\begin{itemize}
    \item Handle unpredictable system behavior
    \item Estimate extreme conditions like overload
\end{itemize}

\textbf{Approach:} Use probabilistic bounds.

\section*{Slide 5: Tail Bounds}
\textbf{Definition:} Tail bounds give an upper bound on probability that a random variable deviates significantly from its expected value.

\textbf{Purpose:}
\begin{itemize}
    \item Quantify risk of extreme events
    \item Assist in system reliability
\end{itemize}

\section*{Slide 6: Types of Tail Bounds}
\begin{enumerate}
    \item Markov’s Inequality
    \item Chebyshev’s Inequality
    \item Chernoff Bound
\end{enumerate}

\section*{Slide 7: Markov’s Inequality}
\[
P(X \geq a) \leq \frac{E[X]}{a}
\]

\textbf{Assumptions:}
\begin{itemize}
    \item $X \geq 0$
    \item $E[X]$ exists
\end{itemize}

\section*{Slide 8: Derivation of Markov’s Inequality}
\[
E[X] = \int_0^\infty P(X \geq t)\, dt
\]

For $t \leq a$:
\[
P(X \geq t) \geq P(X \geq a)
\]

\[
E[X] \geq \int_0^a P(X \geq a)\, dt = a \cdot P(X \geq a)
\]

\[
\Rightarrow P(X \geq a) \leq \frac{E[X]}{a}
\]

\section*{Slide 9: Interpretation of Markov’s Inequality}
\begin{itemize}
    \item Requires only mean
    \item Provides loose bound
    \item Useful with limited information
\end{itemize}

\section*{Slide 10: Chebyshev’s Inequality}
\[
P(|X - \mu| \geq k\sigma) \leq \frac{1}{k^2}
\]

\textbf{Assumptions:}
\begin{itemize}
    \item Mean $\mu$ exists
    \item Variance $\sigma^2$ exists
\end{itemize}

\section*{Slide 11: Derivation of Chebyshev’s Inequality}
Apply Markov to $(X - \mu)^2$:

\[
P((X - \mu)^2 \geq k^2\sigma^2) \leq \frac{E[(X - \mu)^2]}{k^2\sigma^2}
\]

Since:
\[
E[(X - \mu)^2] = \sigma^2
\]

\[
\Rightarrow P(|X - \mu| \geq k\sigma) \leq \frac{1}{k^2}
\]

\section*{Slide 12: Interpretation of Chebyshev’s Inequality}
\begin{itemize}
    \item Works for any distribution
    \item Uses variance
    \item Tighter than Markov
\end{itemize}

\section*{Slide 13: Chernoff Bound}
\[
P(X \geq (1+\delta)\mu) \leq e^{-\frac{\mu \delta^2}{3}}, \quad (0 < \delta < 1)
\]

\textbf{Assumptions:}
\begin{itemize}
    \item Independent random variables
    \item Often Bernoulli trials
\end{itemize}

\section*{Slide 14: Properties of Chernoff Bound}
\begin{itemize}
    \item Exponentially decreasing bounds
    \item Much tighter
    \item Suitable for large-scale systems
\end{itemize}

\section*{Slide 15: Application to Amazon Case Study}
\textbf{Model:}
\begin{itemize}
    \item $X$ = system load (number of users)
\end{itemize}

\textbf{Objective:}
\begin{itemize}
    \item Estimate probability of overload and failure
\end{itemize}

\textbf{Method:} Apply tail bounds.

\section*{Slide 16: System Reliability}
\textbf{Definition:} Probability that a system performs correctly under specified conditions.

\textbf{Use:}
\begin{itemize}
    \item Estimate failure probability
    \item Design safety margins
\end{itemize}

\section*{Slide 17: Key Observations}

\begin{center}
\begin{tabular}{lll}
\toprule
Method & Requirement & Nature \\
\midrule
Markov & Mean only & Loose \\
Chebyshev & Mean + Variance & Moderate \\
Chernoff & Independence & Tight \\
\bottomrule
\end{tabular}
\end{center}

\begin{itemize}
    \item Exact prediction not feasible
    \item Tail bounds provide upper guarantees
\end{itemize}

\section*{Slide 18: Class Participation}
\begin{itemize}
    \item Discussion on tail bounds application
    \item System reliability
    \item Interpretation of probabilistic guarantees
\end{itemize}

\end{document}